\section{Bedürfnisse}
\label{beduerfnisse}
\index{Bedürfnisse}

\subsection{Was sind universelle Bedürfnisse?}

Bedürfnisse sind das, was wir erfüllt brauchen, damit es uns gut geht.

Ein universelles Bedürfnis ist eins, das jeder Mensch kennt~-- auch wenn sich Menschen darin unterscheiden, welche Bedürfnisse sie wie stark erfüllt brauchen.

Echte Bedürfnisse sind nicht an eine konkrete Person gebunden. Es gibt aber durchaus Bedürfnisse, die wir nur mit anderen Menschen zusammen erfüllen können, zum Beispiel unser Bedürfnis nach Gemeinschaft.

Ein Bedürfnis ist nicht an eine konkrete Handlung gebunden. Für jedes Bedürfnis gibt es viele verschiedene Strategien, um sie zu erfüllen~-- und wenn euch nur eine einzige Strategie dafür einfällt, habt ihr das Bedürfnis noch nicht genug verstanden.

\subsection{Liste von Bedürfnissen}
\label{beduerfnisse-liste}

Das ursprüngliche Vokabular stammt von Marshall Rosenberg aus \cite[S.~216~f]{gfk-rosenberg} bzw.~im englischsprachigen Original \cite[S.~210]{nvc-rosenberg}. Das erweiterte Vokabular kommt aus~\cite[S.~75~f]{gfk-dummies}.

\begin{multicols}{2}
    \begin{itemize}
        \item (innerer) Friede
        \item Abwechslung
        \item Achtsamkeit
        \item Akzeptanz
        \item Anerkennung
        \item Anerkennung
        \item Atmen
        \item Aufmerksamkeit
        \item Ausgewogenheit
        \item Ausruhen
        \item Austausch
        \item Authentizität
        \item Balance (von Geben und Nehmen)
        \item Bedeutung
        \item Beitragen
        \item Beteiligung
        \item Bewegung
        \item Bildung
        \item Distanz
        \item Effektivität
        \item Ehrlichkeit/Aufrichtigkeit
        \item Eindeutigkeit
        \item Einklang
        \item Empathie (bekommen)
        \item Engagement
        \item Erfolg (im Sinne von \glqq Gelingen\grqq)
        \item Erholung
        \item Feedback
        \item Feiern (von Erfolgen)
        \item Freiheit
        \item Freude
        \item Freundschaft
        \item Frieden (mit anderen Menschen)
        \item Geborgenheit
        \item Gemeinschaft
        \item Gerechtigkeit
        \item Gesundheit
        \item Gleichbehandlung
        \item Glück
        \item Harmonie
        \item Heilung
        \item Humor
        \item Individualität
        \item Inspiration
        \item Integrität (im Einklang mit den eigenen Werten sein)
        \item Intimität
        \item Klarheit
        \item Kompetenz
        \item Kontakt
        \item Kraft
        \item Kreativität
        \item Körperkontakt
        \item Lebenserhaltung
        \item Leichtigkeit
        \item Lernen
        \item Liebe (erfahren)
        \item Nahrung, Essen
        \item Nähe
        \item Offenheit
        \item Pläne für die Erfüllung der eigenen Träume, Ziele und Werte entwickeln
        \item Respekt
        \item Ruhe
        \item Rückmeldung
        \item Rücksichtnahme
        \item Schlaf
        \item Schutz
        \item Schutz (körperlich)
        \item Schönheit
        \item Selbstbestimmung
        \item Selbstwert
        \item Selbstwirksamkeit
        \item Sexualität
        \item Sexualität
        \item Sicherheit
        \item Sicherheit (körperlich)
        \item Sinn
        \item Spaß
        \item Spiel
        \item Spiel, Spielen
        \item Toleranz
        \item Trauern (über Verluste oder wegen eines Scheiterns)
        \item Trinken
        \item Unterkunft
        \item Unterstützung
        \item Vertrauen
        \item Wertschätzung
        \item Wärme
        \item Zugehörigkeit
        \item Zusammenarbeit
        \item Zärtlichkeit
        \item \glqq Ordnung\grqq{} (im Sinne von Struktur, Klarheit)
        \item die eigenen Träume, Ziele und Werte wählen
        \item persönliches Wachstum
        \item verstanden/gesehen werden
        \item Ästhetik
        \item Übereinstimmung mit den eigenen Werten
    \end{itemize}
\end{multicols}
